\subsection{Intrinsic determinant measures and the complete profile}

Keep the smooth positive probability volume $\mathrm dV$ fixed. For a positive $(n-1,n-1)$-current $T$ of finite mass, let $\Gamma_{0,h}(T)$ denote the determinant measure of \eqref{eq:weak-gaud-functional} computed with a positive auxiliary metric $h$. The current need not be smooth or strictly positive. The following coordinate calculation identifies the measure itself, including on singular coefficient measures.

\Needspace{11\baselineskip}
\begin{lemma}\label{lem:gaud-determinant-measure-intrinsic}\index{Gauduchon duality!intrinsic determinant measure}
For a fixed positive current $T$ and fixed $\mathrm dV$, the measure
$\Gamma_{0,h}(T)$ is independent of the positive auxiliary metric $h$.
This statement includes singular and degenerate currents.
\end{lemma}
\begin{proof}
In a holomorphic frame put
$\lambda=\bigwedge_{j=1}^n(i\,dz_j\wedge d\bar z_j)$,
$h=i\sum H_{ij}\,dz_i\wedge d\bar z_j$ and $\mathrm dV=v\lambda$.
Positive currents have measure coefficients. Choose a positive
Radon measure $\nu$ dominating them. For a smooth real form
$\xi=i\sum\Xi_{ij}\,dz_i\wedge d\bar z_j$, let $Q$ be the
positive semidefinite Hermitian coefficient matrix characterized by
$\xi\wedge T=\operatorname{tr}(\Xi Q)\nu$. Then
\[
 \sigma_h=\operatorname{tr}(HQ)\nu,\qquad
 P_h=\frac{nH^{1/2}QH^{1/2}}{\operatorname{tr}(HQ)}.
\]
The formula holds on $Q\ne0$; zero is assigned when $Q=0$.
Since $h^n=n!\det(H)\lambda$, determinant algebra gives
\begin{equation}\label{eq:intrinsic-gaud-determinant-coordinate}
 \Gamma_{0,h}(T)
 =n(n!)^{-1/n}v^{1/n}(\det Q)^{1/n}\nu.
\end{equation}
This expression has no $H$. Replacing $\nu$ by $L\nu$, where
$L>0$ almost everywhere, replaces $Q$ by $Q/L$, leaving the
measure unchanged; a common dominating measure compares any two choices. The original geometric definition
ensures compatibility under frame changes. Matrix AM--GM bounds it
locally by a constant times a reference trace measure, proving local
finiteness. Compactness gives a finite Radon measure.
\end{proof}
Write $\Gamma_0(T)$ for this intrinsic measure.


\Needspace{18\baselineskip}
\begin{theorem}\label{thm:constant-mass-entire-determinant-profile}\index{Gauduchon duality!exact density-root profile}
Let $\beta$ be smooth positive Hermitian with
$\ddc(\beta^j)=0$ for $1\leq j\leq n-1$,
$V=\int_X\beta^n$, and let $h=\beta+\ddc u>0$ for a smooth
global real function $u$. Then
\[
 \mathscr G_h=\mathscr G_\beta,\qquad
 \{\Gamma_{0,h}(T):T\in\mathscr G_h\}
 =\{\Gamma_{0,\beta}(T):T\in\mathscr G_\beta\}.
\]
Every measure in this common family has a density $k_T$ satisfying
\[
 k_T\geq0,\qquad
 \|k_T\|_{L^{n/(n-1)}(\mathrm dV)}\leq V^{-1/n}.
\]
For every real or complex $p\in L^n(\mathrm dV)$,
\begin{equation}\label{eq:constant-mass-determinant-profile}
 \sup_{T\in\mathscr G_h}\int_X|p|\,\Gamma_0(T)
       =V^{-1/n}\|p\|_{L^n(\mathrm dV)}.
\end{equation}
All suprema retain singular and degenerate tests.
\end{theorem}
\begin{proof}
Distributional integration by parts gives
$\int h\wedge T=\int\beta\wedge T+\int u\ddc T
=\int\beta\wedge T$. This proves equality of the actual current
sets; Lemma~\ref{lem:gaud-determinant-measure-intrinsic} proves equality of their measure sets.

Take a smooth strictly positive real $r$ and put $M=\int r^n \mathrm dV$.
By \cref{thm:tw-smooth} there are a smooth potential $w$, a positive form $A=\beta+\ddc w$ and a positive coefficient in its prescribed-volume equation. The mass-invariance hypothesis and \cref{thm:mass-invariance-characterization} determine that coefficient, giving
\begin{equation}\label{eq:constant-mass-reference-equation}
 A^n=\frac{Vr^n}{M}\,\mathrm dV.
\end{equation}
In particular, the total mass of $A$ is $V$. No estimates uniform in a degenerating smooth density are required.

For every $T\in\mathscr G_\beta$, integration by parts gives
$\int A\wedge T=1$. In an $A$-unitary frame, matrix AM--GM gives
$\det(P_A)^{1/n}\leq1$, including singular coefficient densities.
Thus as Radon measures
\[
 \Gamma_0(T)\leq
       (\mathrm dV/A^n)^{1/n}(A\wedge T).
\]
Multiplying by $r$ and using \eqref{eq:constant-mass-reference-equation} yields
\begin{equation}\label{eq:constant-mass-determinant-dual}
 \int r\,\Gamma_0(T)\leq V^{-1/n}\|r\|_{L^n(\mathrm dV)}.
\end{equation}
Replacing a smooth nonnegative $r$ by $r+\eta$ and letting
$\eta\downarrow0$ extends this bound to nonnegative smooth tests.
For smooth real or complex $s$, smooth positive approximations
$\sqrt{|s|^2+\eta^2}$ imply
$|\int s\,\Gamma_0(T)|\leq V^{-1/n}\|s\|_n$.
Consequently integration against $\Gamma_0(T)$ extends continuously
from smooth functions to $L^n$. Its positive dual representative is
$k_T\in L^{n/(n-1)}$ with the stated bound. Equality on smooth
tests identifies the two Radon measures, proving absolute continuity.
The integrals of arbitrary $L^n$ roots are therefore well-defined.

For the reverse inequality with $r>0$ smooth, let
$T_r=A^{n-1}/V$. It is smooth and positive. Write $E=\ddc w$. For $j\leq n-1$, expansion and closedness of $E$ give
\[
 \ddc(A^j)=\sum_{l=0}^j\binom jl\ddc(\beta^{j-l})\wedge E^l=0.
\]
The term with $j=l$ differentiates the constant one. Thus $\ddc T_r=0$, and smooth integration by parts gives
\[
 \int_X\beta\wedge A^{n-1}
 =\int_X A^n-\int_X\ddc w\wedge A^{n-1}=V.
\]
Consequently $T_r\in\mathscr G_\beta$.
Its positive root is $\gamma=V^{-1/(n-1)}A$, so
\[
 k_{T_r}=(\gamma^n/\mathrm dV)^{(n-1)/n}
       =V^{-1/n}\frac{r^{n-1}}{\|r\|_n^{n-1}}.
\]
It attains equality in \eqref{eq:constant-mass-determinant-dual}. For a nonnegative
$p\in L^n$, take smooth positive $r_j\to p$ in $L^n$. The common
dual bound controls $\int(p-r_j)k_{T_{r_j}}\mathrm dV$, proving the lower
bound in \eqref{eq:constant-mass-determinant-profile}; the upper bound already holds for
every current. Apply the same argument to $|p|$ for complex roots.
\end{proof}
